\scanpage{281}
envisagé où le schéma de Néron $\uuline{\operatorname{N\acute{e}r}}_{X/S}$ existe, le foncteur $\uline{\operatorname{N\acute{e}r}}^{+\prime\prime}_{X/S}$ est associé au sous-schéma ouvert induit $\uuline{\operatorname{N\acute{e}r}}^{+}_{X/S}$ de $\uuline{\operatorname{N\acute{e}r}}_{X/S}$, image de $\uuline{\operatorname{Pic}}^{+}_{X/S}$. Un élément de $\uline{\operatorname{N\acute{e}r}}^{+\prime\prime}_{X/S}(S)=\operatorname{N\acute{e}r}^{+\prime\prime}(X/S)$ est appelé une \uline{polarisation} \deleted{\unsure{positive}} de $X/S$. Elles \inserted{forment un pré-monoïde (sans élément neutre) commutatif}\inserted{\unsure{plongé dans le groupe} $\uline{\operatorname{N\acute{e}r}}''(X/S)$}.

Deux polarisations $\xi,\xi'$ sont dites \uline{équivalentes} s'il existe des entiers $n,n'>0$ tels que $n\xi=n'\xi'$. Les \inserted{classes de} polarisations correspondent donc aux demi-droites d'un cône convexe $C$ \deleted{[ouvert]} dans un espace vectoriel $V$ sur $\unsure{\mathbb Q|\mathbf Q}$ engendré par $C$, dont les pts de $C$ sont ``\unsure{intérieurs}'' \unsure{---}.

Polarisation \uline{très ample} (pour une polarisation donnée, il en existe tjrs une équivalente qui est \uline{très ample})
\figure{\begin{tikzpicture}[x=1cm,y=1cm]
\draw (0,0.4)..controls (0.1,0.8) and (0.4,0.95)..(0.65,1.05); \draw (0.65,1.05)..controls (0.4,1.1) and (0.4,1.25)..(0.35,1.5);
\draw (2,0)--(2,0.3)..controls (2.0,0.48) and (2.2,0.52)..(2.45,0.52)--(3.45,0.55);
\end{tikzpicture}}
\mspage{280}
